% Copyright 2014 Daniel Nikpayuk
\documentclass[twoside]{article}
  \setlength{\oddsidemargin}{0in}\setlength{\evensidemargin}{0in}
  \setlength{\topmargin}{0in}
  \setlength{\headheight}{0in}\setlength{\headsep}{0in}
  \setlength{\textwidth}{6.5in}\setlength{\textheight}{9in}
\usepackage{amssymb}
\usepackage{graphicx}
\usepackage{hyperref}
\title{Summation}
\author{Daniel Nikpayuk}
\date{December 8, 2014}
\begin{document}
\maketitle

\begin{figure}[h]
\centering
\includegraphics[width=1in]{cc-by-nc.png}\\[0.1in]
\tiny This article is licensed under \\
\href{http://creativecommons.org/licenses/by-nc/4.0/}
{Creative Commons Attribution-NonCommercial 4.0 International.}\\[0.3in]
\end{figure}

Summation for the many who are not well acquainted with it, as 
old friends are with each other, do not much enjoy the sight of 
the, as some might say, horrendous notation given to it.  I'm 
talking about the notorious \emph{sigma} notation $  \Sigma  $, 
which is meant \mbox{to---if} learned and used 
\mbox{correctly---make} handling ugly sums a walk in the park.

Such people often invoke the ``path of least resistence'' 
arguement by which they say there is no need to introduce new 
notation maintaining the existing ellipsis notation, i.e. 
$  a_1+a_2+\ldots +a_n  $, works just fine since it is less 
mechanical, easier to manipulate, and therefore faster to use.  
The point being: if it's not broke, and it gets the job done,
why fix it?

It is not the author's intention to turn this observation into 
a debate.  Instead, the intention of this article is to make 
refinements in understanding the notation and why it exists 
(other than as a shortform of the horizontally longer ellipsis 
representation), as well as demonstrating the power it holds 
(most notably with finding closed forms for multiple sums).

\section{Sigma notation}

We will be working with what is called the generalized sigma
notation.  It turns out that there are two varieties which
are worth introducing since they will both help a great deal in
manipulating sums.  The first provides us with a
set-theoretic perspective, and the second handles the notation 
with the idea of properties. 

So to begin, the following notation: 
$$  \sum_{k\in I}a_k  $$ is simply an abbreviation for the sum 
of all terms $  a_k  $ such that $  k  $ is an \emph{integer} 
that belongs $  \in  $ to the set $  I  $, where $  I  $ is 
called the \emph{index set} and is a finite subset of the 
integers $  \mathbb{Z}  $.  Just to clarify, the $  \Sigma  $ part of 
the sigma notation is called the \emph{summand}.

When forming the index set, one may simply collect integers
together, even if it's for no apparent reason.  But
one can also collect the necessary integers together in a
systematic way by means of some known property $  P(k)  $.
In that case, as in the book ``Concrete Mathematics''
for which this article is inspired (and is otherwise referenced
as the bracketed index \cite{gkp}), we write:
$$  \sum_{P(k)}a_k   $$ and add terms $  a_k  $ to our sum 
whenever $  k  $ is an integer satisfying the given property 
$  P(k)  $.  The mindful reader at this point might ask what 
happens when no integers satisfy the property $  P(k)  $?  
Well, since we aren't adding anything to the sum, we will 
again borrow from \cite{gkp} and define such a structure to 
be \emph{zero}.  The same thing happens when we sum terms over 
an empty index set, so we should name this peculiarity the 
\emph{empty sum}.

For concreteness, the following example is provided.  Let 
$  a_k=1  $ and $  I=\{ -2,-1,0,1,2,3\}  $, then
$$  \sum_{k\in I}1=1+1+1+1+1+1=\sum_{-2\le k\le 3}\!\!\!\! 1  $$

We now take this idea a bit further to introduce double sums.  
These types of sums are actually pretty simple, they arise 
when we consider the possibility of the term of a series
$  \sum_{j\in H}b_j  $ being equal to a sum itself.  With this 
comes a new consideration though, the term $  b_j  $ as a sum
will have its own term $  a_{j,k}  $ and its own index set 
$  I  $, both of which depend on $  k  $, but \emph{can} also 
depend on $  j  $.
$$  \sum_{j\in H}\sum_{k\in I_j}\!a_{j,k}
\mbox{\ \ \ or\ \ \ }\sum_{Q(j)}\sum_{P(j,k)}\!\! a_{j,k}  $$

Noting one more thing before we move on, I will actually use 
the better known delimited form of the sigma notation in this 
article, ie. 
$$  \sum_{k=m}^na_k=a_m+a_{m+1}+a_{m+2}+\ldots +a_{n-1}+a_n  $$
but only to \emph{display} equations in 
their finished form.

\subsection{Manipulation}

The goal of our notation now that we have it, is to be
able to manipulate the complicated structure of summation
in a concise and easy way.  With this in mind, let's see
what we can do.   

We begin actually with double sums.  The essential thing to 
realize here, is that a double sum is a sum of sums, which 
in the end is just another sum (say that ten times fast).  
Once this is thought out, one would expect that there is a 
reasonable \emph{single summand} representation of certain 
(infact many) double sums.  This of course is true.  

We approach this idea, of a single summand representation of 
a double sum, by extending the original definitions, by which 
we extend the term $  a_j  $ to $  a_{j,k}  $ and the index set 
$  I  $ of \emph{integers} to an index set $  I  $ of 
\emph{ordered pairs of integers} (or a property $  P(k)  $ to a 
property $  P(j,k)  $).  Here, an example of this alternate 
approach would be of some help.  How about letting 
$  a_{j,k}=j+k  $ and
$$  I=\{ (0,0),(0,1),(0,2),(1,1),(1,2),(2,2)\}  $$
so then 
$$  \sum_{(j,k)\in I}\!\! (j+k)=\sum_{0\le j\le k\le 2}\!\!\!\!\! (j+k)  $$
$$  =(0+0)+(0+1)+(0+2)+(1+1)+(1+2)+(2+2)  $$
See!  easy as $  \pi  $.   

With this extension in hand, we are led to: 

\begin{eqnarray*}
\sum_{(m\le j\le k\le n)}\!\!\!\!\!\!\!\!\! a_{j,k}                         
 & = & \sum_{(m\le k\le n)}\sum_{(m\le j\le k)}\!\!\!\!\! a_{j,k} \\
 & = & \sum_{(m\le j\le n)}\sum_{(j\le k\le n)}\!\!\!\! a_{j,k}
\end{eqnarray*} for some integer constants $  m  $ and $  n  $.   
This handy little 3-way equation is called the \emph{summand 
switch law}.  There are infact broader summand switch laws, 
but for our applications, the one presented here is sufficient. 

The truth of this law isn't exactly ``obvious,'' so it helps if
you look at it this way: since all three formulae are sums of the 
terms $  a_{j,k}  $, we only have to show that whenever 
$  a_{j,k}  $ belongs to any of the above sums for particular 
values $  j,k  $ they also belong to the two others (the same 
number of times of course). 

We can actually ``streamline'' this proof, not by showing that 
for each sum containing any given term the other two sums 
contain that given term, but by showing that the properties 
cummulated within each sum, for each sum (that need to be 
satisfied in order for the term to belong), are equivalent.  
Let's see, how do we do this?\ldots Ah, that's easy!
It is first worthwhile to note that in general 
$$  (a\le b\le c)  $$ is shortform for
$$  (a\le b)\mbox{ \emph{and} } (b\le c)  $$
With this we observe that
$$  (m\le k\le n)\mbox{ \emph{and} } (m\le j\le k)  $$ 
is true if and only if
$$  (m\le j)\mbox{ \emph{and} } (j\le k)\mbox{ \emph{and} } (k\le n)  $$ 
is true, or equivalently 
$$  (m\le j\le k\le n)  $$ which working backwards is the same as
$$  (m\le j\le n)\mbox{ \emph{and} } (j\le k\le n)  $$
A symmetric arguement may be given in the case of the properties
being false to show their complete, and total logical equivalence.  

Applying this now proven law to our above example we see that

\begin{eqnarray*}
 &   & \sum_{0\le j\le k\le 2}\!\!\!\!\! (j+k)                  \\
 & = & \sum_{(0\le k\le 2)}\sum_{(0\le j\le k)}\!\!\!\! (j+k)   \\ 
 & = & [\ (0+0)\ ]\ +\ [\ (0+1)+(1+1)\ ]                        \\
 &   & +\ [\ (0+2)+(1+2)+(2+2)\ ]                               \\
 & = & [\ (0+0)+(0+1)+(0+2)\ ]                                  \\
 &   & +\ [\ (1+1)+(1+2)\ ]\ +\ [\ (2+2)\ ]                     \\ 
 & = & \sum_{(0\le j\le 2)}\sum_{(j\le k\le 2)}\!\!\!\! (j+k)             
\end{eqnarray*} 

Moving on with our manipulation laws, it turns out that there are 
three major ones to know when looking at single summand sums.  
The first two are straight-forward, so I won't spend time with any 
explanations, they work out to be:

\begin{eqnarray}
\sum_{k\in I}(a_k+b_k) & = & \sum_{k\in I}a_k+\sum_{k\in I}b_k \\
\sum_{k\in I}ca_k & = & c\sum_{k\in I}a_k 
\end{eqnarray} Both of these equations have names, $  (1)  $ is 
called the \emph{associative law}, and $  (2)  $ is called the 
\emph{distributive law}.  The use of the phrase ``single summand 
sums'' (above) should be noted since it implies that the sums in 
these equations may infact be double sums in disguise.  This works 
as long as we consider the possibility of the index variable 
$  k  $ being shortform for $  k=(k_1,k_2)  $.  The truth of these 
extended laws may be verified with the set theoretic arguements 
applied to the index set.

Hmmm\ldots what could it possibly mean if I said there were
three laws and only included two as being straight-forward?
Yes! that's right, the third law, called the 
\emph{commutative law}, will need some clarification.  
It takes the form:

\begin{eqnarray}
\sum_{k\in I}a_k & = & \sum_{p(k)\in I}\!\!\! a_{p(k)}
\end{eqnarray}Wait a minute though, isn't this just
substituting one variable for another?  Nope! consider:
$$  \sum_{0\le k\le 5}\!\!\! a_k\stackrel{?}{=}
\sum_{0\le k^2\le 5}\!\!\!\! a_{k^2}  $$
the former (on the left) equals $$  a_0+a_1+a_2+a_3+a_4+a_5  $$
while the latter (on the right) is equal to $$  a_{4}+a_{1}+a_0+a_1+a_4  $$
The point is that not just any substitution will do,
there needs to be a one-to-one correspondence, i.e. for
every integer (or pair of integers) $  k_1  $ belonging to 
the index set of summation, there is \emph{exactly} one 
integer (or pair) $  k_2  $ such that $  k_1=p(k_2)  $, 
otherwise you could add the same term more than once or 
not at all.  A good example of a ``proper'' substitution, 
is setting $  p(k)=n-k  $ for the sum
$$  \sum_{0\le k\le n}\!\!\! a_k  $$ 
since for each $  k_1\in I  $ (the index set)
there exists only one $  k_2  $ such that $  k_1=n-k_2  $,
or equivalently, there is only one corresponding
$  n-k_2\in I  $. Substituting this in, we have
$$  \sum_{0\le k_1\le n}\!\!\!\! a_{k_1}
 =\sum_{0\le n-k_2\le n}\!\!\!\!\!\!\! a_{n-k_2}  $$
So to restate the obvious, the term $  a_{n-k_2}  $ is 
added to the sum if and only if 
$$  (0\le n-k_2\le n)\mbox{ is true, }  $$ 
now looking more closely at this particular property 
$$  (0\le n-k_2\le n)  $$ we see that
$$  (0\le n-k_2\le n)\Longleftrightarrow (0\le k_2\le n)  $$
which follows from again using the fact that
$$  (0\le n-k_2\le n)  $$ is shortform for
$$  (0\le n-k_2)\mbox{ \emph{and} }(n-k_2\le n)  $$
but also
$$  (0\le n-k_2)\Longleftrightarrow (k_2\le n)  $$
$$  (n-k_2\le n)\Longleftrightarrow (0\le k_2)  $$ 
What all of this then means is that
$$  \sum_{0\le k_1\le n}\!\!\!\! a_{k_1}
=\sum_{0\le k_2\le n}\!\!\!\! a_{n-k_2}  $$ or equivalently:
$$  \sum_{0\le k\le n}\!\!\! a_{k}
=\sum_{0\le k\le n}\!\!\! a_{n-k}  $$

The important thing to remember in this chain of reasoning is 
the \emph{method of derivation} rather than the specific result 
itself, because with it you can derive many other specific 
commutative laws, for example setting \mbox{$  p(k)=k+c  $} will 
always work, as well as using the extensions 
\mbox{$  p(j,k)=(j+c,k+c)  $} or \mbox{$  p(j,k)=(n-j,n-k)  $} 
(for some integer constants $  c  $ and $  n  $).  Having this 
said, I will from now on entirely skip any other similar 
derivation when proving an equation.

\section{Practice}

We are finally ready to demonstrate the value of our now
furnished sigma notation.  It should be noted that the two
famous series
$$  \sum_{1\le k\le n}\!\!\! k=\frac{n(n+1)}{2}\mbox{\ \ \ and\ \ }  
\sum_{0\le k\le n}\!\!\! x^k=\frac{1-x^{n+1}}{1-x^{}\ \ \ \ }  $$
are taken for granted since proofs can be supplied very
easily from very many sources.

Let's start off nice and simple with
$$  \sum_{0\le j\le k\le n}\!\!\!\!\! 1  $$ 

\begin{eqnarray*}
 & = & \sum_{(0\le k\le n)}\sum_{(0\le j\le k)}\!\!\! 1              \\
 & = & \sum_{(0\le k\le n)}\!\!\!\! (k+1)                            \\
 & = & \sum_{(0\le k\le n)}\!\!\!\! k+\sum_{(0\le k\le n)}\!\!\!\! 1 \\
 & = & \frac{n(n+1)}{2}+(n+1)                      \\
 & = & \frac{(n+1)(n+2)}{2} 
\end{eqnarray*}

As our second example, let us find a closed form for
$$  \sum_{0\le j\le k\le n}\!\!\!\!\! (j+k)  $$
We could solve this old friend of ours\footnote{As far as 
articles are concerned.} one sum at a time (as in the 
previous example), but there is a more clever method.    

We start with the permutation on $  \mathbb{Z}^2  $
$$  p(j,k)=(n-j,n-k)  $$ which after implementing the 
extended commutative law gives us

\begin{eqnarray*}
 &   & \sum_{0\le n-j\le n-k\le n}\!\!\!\!\!\!\!\!\!\!\! [(n-j)+(n-k)]  \\
 & = & \sum_{0\le k\le j\le n}\!\!\!\!\! [2n-(j+k)]                     \\
 & = & 2n\!\!\!\!\!\!\sum_{0\le k\le j\le n}\!\!\!\!\!\! 1
-\sum_{0\le k\le j\le n}\!\!\!\!\! (j+k) 
\end{eqnarray*} From here we observe, for the purpose of 
substitution, that

\begin{eqnarray*}
\sum_{0\le k\le j\le n}\!\!\!\!\! (j+k) 
 & = & \sum_{0\le k\le j\le n}\!\!\!\!\! (k+j) \\
 & = & \sum_{0\le j\le k\le n}\!\!\!\!\! (j+k)
\end{eqnarray*} leading us (with a little bit of old-fashioned
 algebra) to

\begin{eqnarray*}
\sum_{0\le j\le k\le n}\!\!\!\!\! (j+k) 
 & = & n\!\!\!\!\!\!\sum_{0\le k\le j\le n}\!\!\!\!\! 1 \\
 & = & \frac{n(n+1)(n+2)}{2}
\end{eqnarray*} Interestingly, this solution may be applied 
towards the well known series $  \sum_{k=1}^nk^2  $:

\begin{eqnarray*}
 &   & \frac{n(n+1)(n+2)}{2}                                           \\
 & = & \sum_{0\le j\le k\le n}\!\!\!\!\! (j+k)                         \\
 & = & \sum_{(0\le k\le n)}\sum_{(0\le j\le k)}\!\!\!\! (j+k)          \\
 & = & \sum_{(0\le k\le n)}\sum_{(0\le j\le k)}\!\!\!\! j               
+\sum_{(0\le k\le n)}\sum_{(0\le j\le k)}\!\!\!\! k                    \\
 & = & \sum_{(0\le k\le n)}\!\!\!\!\!\frac{k(k+1)}{2}                   
+\sum_{(0\le k\le n)}\!\!\!\!\! k(k+1)                                 \\
 & = & \frac{3}{2}\sum_{(0\le k\le n)}\!\!\!\! (k^2+k)                 \\
 & = & \frac{3}{2}\sum_{(0\le k\le n)}\!\!\!\!\! k^2+\frac{3n(n+1)}{4}
\end{eqnarray*} 
$$  \Longrightarrow \sum_{k=1}^nk^2=\frac{n(n+1)(2n+1)}{6}  $$

For our next example, let's try something a little  
different: 
$$  S_n=\sum_{0\le k\le n}\!\! (k+1)x^k  $$
A good start is to relegate and replace this series with its alter ego:
$$  \sum_{0\le j\le k\le n}\!\!\!\!\! x^k  $$ this claim holding of
course because 
$$  S_n=\sum_{(0\le k\le n)}\sum_{(0\le j\le k)}\!\!\!\! x^k  $$

Venturing towards something new, we will partition this sum into 
two sums in two ways (why you ask?\mbox{\ldots be patient)}.  
We do this by partitioning the terms of the sum, effectively 
dividing the sum into two.  The first division occurs when we use 
the fact that the index variable $  k  $ of the general term 
$  x^k  $ will either be $  k=n  $ or $  k<n  $:

\begin{eqnarray*}
\sum_{0\le j\le k\le n}\!\!\!\!\! x^k             
 & = & \sum_{0\le j\le k=n}\!\!\!\!\! x^n        
 +\sum_{0\le j\le k<n}\!\!\!\!\! x^k            \\
 & = & x^n\!\!\!\sum_{0\le j\le n}\!\!\! 1             
 +\sum_{0\le j\le k\le n-1}\!\!\!\!\!\!\!\! x^k \\
\end{eqnarray*} the second way to partition the terms is to realize 
that the index variable $  j  $ will either be $  0=j  $ or $  0<j  $ 
so that 

\begin{eqnarray*}
\sum_{0\le j\le k\le n}\!\!\!\!\! x^k      
 & = & \sum_{0=j\le k\le n}\!\!\!\!\! x^k   
 +\sum_{0<j\le k\le n}\!\!\!\!\! x^k      \\
 & = & \sum_{0\le k\le n}\!\!\! x^k        
 +\sum_{1\le j\le k\le n}\!\!\!\!\! x^k    
\end{eqnarray*} We then convert
$$  \sum_{1\le j\le k\le n}\!\!\!\!\! x^k
\mbox{\ \ \ to\ \ \ }\sum_{0\le j\le k\le n-1}\!\!\!\!\!\!\!\! x^k  $$ 
with $$  p(j,k)=(j+1,k+1)  $$ which is a permutation on $  \mathbb{Z}^2  $.
This is carried out as follows:

\begin{eqnarray*}
\sum_{1\le j\le k\le n}\!\!\!\!\! x^k 
 & = & \sum_{1\le j+1\le k+1\le n}\!\!\!\!\!\!\!\!\!\!\! x^{k+1}        \\
 & = & x\!\!\!\!\!\!\!\!\!\sum_{0\le j\le k\le n-1}\!\!\!\!\!\!\!\! x^k
\end{eqnarray*} which after substituting back in and equating
the two derived sides, leaves us with
$$  x^n\sum_{j=0}^n1+S_{n-1}=\sum_{k=0}^nx^k+xS_{n-1}  $$
reducing to
$$  \sum_{k=0}^n(k+1)x^k=\frac{1-x^{n+2}}{(1-x)^2}
-\frac{(n+2)x^{n+1}}{1-x^{}}  $$

As our last example, we have:
$$  T_n=\sum_{0\le j\le k\le n}\!\!\!\!\! x^{j+k}  $$
which will tell us as to whether or not our  adventurous 
partitioning method of the last example was just luck or if it has 
some worth.  Repeating the process provides us with the equation
$$  x^n\sum_{j=0}^nx^j+T_{n-1}=\sum_{k=0}^nx^k+x^2T_{n-1}  $$
so it appears this method is of use after all.
$$  \sum_{0\le j\le k\le n}\!\!\!\!\! x^{j+k}
=\frac{1-x^{n+1}}{1-x^{}\ \ \ \ }\cdot\frac{1-x^{n+2}}{1-x^2\ \ \ }  $$

\newpage

\section{Challenge Problems}
Find closed forms for the following:

\begin{enumerate}
\item $$  \sum_{k=-3}^n[(k+4)^2+(k-2)+3]  $$
\item $$  \sum_{k=2}^n(x+1)^{k+1}  $$
\item $$  \underbrace{\sqrt{2\sqrt{2\sqrt{2\sqrt{\ldots\sqrt{2}}}}}}
_{n^{\#}\mbox{ of times}}  $$
\item $$  \sum_{k=1}^nkx^k  $$
\item $$  \sum_{k=1}^nk^3  $$
      hint: Commute with $  p(k)=n-k  $, and expand the term.
\item $$  \sum_{0\le j,k\le n}\!\!\!\! 1  $$
      hint: $  (0\le j,k\le n)  $

      $  \Longleftrightarrow (0\le j\le n) $ 
      \emph{and} $  (0\le k\le n)  $ 
\item $$  \sum_{0\le j\le k\le\ell\le n}\!\!\!\!\!\!\!\! x^{j+k+\ell}  $$
      hint: Use the partitioning method.
\item If $  H_n  $ is notation such that
$$  H_n=\sum_{k=1}^n\frac{1}{k}  $$ then prove the following: 
$$  \sum_{n=1}^mH_n=(m+1)H_m-m  $$
      hint: Apply the summand switch law.
\end{enumerate}

\begin{thebibliography}{99}
\bibitem{gkp} R.L. Graham, D.E. Knuth, O. Patashnik.  Concrete 
         Mathematics.  Addison-Wesley Publishing (1994).
\end{thebibliography}

\end{document}

